% Hausdorff Measure --- Evans & Gariepy, Chapter 2 (pp. 60--78)
% Lecture notes style, Fall 2024, Mon/Wed schedule
% \input'd by main.tex

\section{September 16, 2024}\label{sec:lec6}

Brian asked me over lunch what Hausdorff measure even \emph{is}, and my
attempt to explain it made me realize I did not understand the definition as
well as I thought. Specifically, I could not say why the limit $\delta \to 0$ matters. Humbling.

Here is the problem. Lebesgue measure only ``sees'' $n$-dimensional things in
$\R^n$. A curve in $\R^3$? Lebesgue measure says it has zero volume and walks away. Same for the surface of a
sphere or a fractal. You need something that can handle non-integer
dimensions. That is Hausdorff measure.

\subsection{Hausdorff pre-measure and Hausdorff measure}%
\label{subsec:hausdorff-def}

We build $\mathcal{H}^s$ in two steps: approximate at scale $\delta$,
then take a limit.

\begin{definition}[Hausdorff pre-measure]\label{def:hausdorff-premeasure}
Let $A \subseteq \R^n$, $0 \le s < \infty$, $0 < \delta \le \infty$.
Define
\[
  \mathcal{H}^s_\delta(A)
  = \inf\!\left\{\, \sum_{j=1}^\infty \alpha(s)
    \left(\frac{\diam C_j}{2}\right)^{\!s} :
    A \subseteq \bigcup_{j=1}^\infty C_j,\;
    \diam C_j \le \delta \,\right\},
\]
where $\alpha(s) = \pi^{s/2}/\Gamma(s/2+1)$ is the normalization constant
chosen so that $\mathcal{H}^n$ agrees with $\mathcal{L}^n$.
\end{definition}

I keep forgetting the normalization. Writing it out one more time:
$\alpha(n) = \pi^{n/2}/\Gamma(n/2+1)$. For integer $n$ this is the volume
of the unit ball in $\R^n$. So $\alpha(1) = 2$, $\alpha(2) = \pi$,
$\alpha(3) = 4\pi/3$. I have a flashcard for this now.

\begin{definition}[$s$-dimensional Hausdorff measure]%
\label{def:hausdorff-measure}
\[
  \mathcal{H}^s(A)
  = \lim_{\delta \to 0^+} \mathcal{H}^s_\delta(A)
  = \sup_{\delta > 0} \mathcal{H}^s_\delta(A).
\]
The limit exists (possibly $= +\infty$) because $\mathcal{H}^s_\delta(A)$
is non-decreasing as $\delta \to 0$: shrinking $\delta$ restricts the
admissible covers.
\end{definition}

As $\delta$ shrinks, covers must use smaller sets, which can
only make the infimum larger. Why force small covers at all? Because at large scales, covering sets can ``cheat'' by straddling geometric features that only become visible at finer resolution. But I am not sure I could make that precise without more work.

\begin{remark}\label{rem:hausdorff-sanity}
Some sanity checks:
\begin{itemize}[nosep]
  \item $\mathcal{H}^0$ is counting measure.
        Each nonempty covering set contributes
        $\alpha(0)(\diam C_j / 2)^0 = 1$.
  \item $\mathcal{H}^1$ on $\R^1$ agrees with $\mathcal{L}^1$.
        Not hard but takes a small argument --- intervals are the most
        efficient covers on the real line.
  \item $\mathcal{H}^n$ on $\R^n$ agrees with $\mathcal{L}^n$.
        That's the content of the isodiametric inequality below.
  \item $\mathcal{H}^s \equiv 0$ on $\R^n$ for $s > n$.
        Cover by cubes of side $\delta$; each has diameter
        $\delta\sqrt{n}$. The sum goes to $0$ as $\delta \to 0$ because
        the exponent $s > n$ beats the count.
\end{itemize}
\end{remark}

$\mathcal{H}^s$ is a Borel regular outer measure on $\R^n$. Countable
subadditivity is straightforward. Borel regularity takes more work
(Evans, p.\ 62).


\subsection{Hausdorff dimension}\label{subsec:hausdorff-dim}

\begin{lemma}[Dimension threshold]\label{lem:dimension-threshold}
Let $A \subseteq \R^n$.
\begin{enumerate}[nosep]
  \item If $\mathcal{H}^s(A) < \infty$, then $\mathcal{H}^t(A) = 0$
        for all $t > s$.
  \item If $\mathcal{H}^t(A) > 0$, then $\mathcal{H}^s(A) = +\infty$
        for all $0 \le s < t$.
\end{enumerate}
\end{lemma}

\begin{proof}[Proof sketch]
For (i): let $\{C_j\}$ be a $\delta$-cover of $A$. Then
\[
  \sum_j \alpha(t)\!\left(\frac{\diam C_j}{2}\right)^{\!t}
  \le \left(\frac{\delta}{2}\right)^{\!t-s}
  \frac{\alpha(t)}{\alpha(s)}
  \sum_j \alpha(s)\!\left(\frac{\diam C_j}{2}\right)^{\!s}.
\]
Take the infimum: $\mathcal{H}^t_\delta(A) \le C \cdot \delta^{t-s}
\cdot \mathcal{H}^s_\delta(A)$. Send $\delta \to 0$. The
$\delta^{t-s}$ factor kills everything since $t - s > 0$ and
$\mathcal{H}^s(A) < \infty$. Part (ii) is the contrapositive. \qedhere
\end{proof}

So there is a critical value of $s$ where $\mathcal{H}^s$ jumps from
$+\infty$ to $0$. A phase transition, if you want to be dramatic about it.

\begin{definition}[Hausdorff dimension]\label{def:hausdorff-dim}
\[
  \dim_{\mathcal{H}}(A)
  = \inf\!\bigl\{ s \ge 0 : \mathcal{H}^s(A) = 0 \bigr\}
  = \sup\!\bigl\{ s \ge 0 : \mathcal{H}^s(A) = +\infty \bigr\}.
\]
\end{definition}

At the critical dimension itself, $\mathcal{H}^s(A)$ can be $0$, $\infty$,
or anything in between. The middle-thirds Cantor set has
$\dim_\mathcal{H} = \log 2 / \log 3$ and its Hausdorff measure at that
dimension is $1$. That is special. Most sets are not so well-behaved at the critical exponent.

Examples:
\begin{itemize}[nosep]
  \item Smooth $k$-dimensional submanifolds of $\R^n$:
        $\dim_\mathcal{H} = k$.
  \item Middle-thirds Cantor set:
        $\dim_\mathcal{H} = \log 2 / \log 3 \approx 0.631$.
  \item Countable sets: $\dim_\mathcal{H} = 0$.
  \item $\R^n$: $\dim_\mathcal{H} = n$. Obviously.
\end{itemize}

I tried computing $\mathcal{H}^{\log 2/\log 3}$ of the Cantor set
directly but got tangled in the covering argument. The upper bound is
easy (use the $2^k$ remaining intervals at level $k$). The lower bound
needs more care: you have to show those ``natural'' covers are actually
optimal, and that is where I got stuck. I think the density argument in Rogers's book (Chapter~1, around Theorem~30) handles it, but I have not worked through it line by line.

TODO: work through the Cantor set dimension calculation carefully. Check
Rogers's book for the density argument.


% ===========================================================================
\section{September 18, 2024}\label{sec:lec7}

The isodiametric inequality via Steiner
symmetrization is one of the cleanest proofs I have seen all year. I was sitting in the back row of KAP 414 with my laptop half-closed and I put it away to watch Evans write it out on the board.

\subsection{Steiner symmetrization}\label{subsec:steiner}

\begin{definition}[Steiner symmetrization]\label{def:steiner}
Let $A \subseteq \R^n$ be Borel and $a \in S^{n-1}$ a unit vector. The
\emph{Steiner symmetrization} of $A$ with respect to the hyperplane
$\{x : x \cdot a = 0\}$ is obtained by replacing, for each line parallel
to $a$, the intersection of that line with $A$ by a symmetric interval of
the same length centered on the hyperplane.

Precisely: decompose $x = y + ta$ where $y \perp a$, $t \in \R$. For
each $y$, let $A_y = \{t : y + ta \in A\}$. Then
\[
  S_a(A) = \bigl\{\, y + ta : y \perp a,\;
    |t| \le \tfrac{1}{2}\mathcal{L}^1(A_y) \,\bigr\}.
\]
\end{definition}

Picture it: slide each vertical ``slice'' of $A$ so it is centered on the
hyperplane. Each slice keeps its length, so volume is preserved.
But the set gets rounder. The diameter can only shrink or stay the same.

\begin{lemma}\label{lem:steiner-props}
Let $A \subseteq \R^n$ be Borel.
\begin{enumerate}[nosep]
  \item $\diam\, S_a(A) \le \diam\, A$.
  \item $\mathcal{L}^n(S_a(A)) = \mathcal{L}^n(A)$.
\end{enumerate}
\end{lemma}

\begin{proof}[Proof sketch]
Part (ii) is immediate from Fubini: each slice $A_y$ gets replaced by an
interval of the same $\mathcal{L}^1$-measure. Integrate over $y$.

Part (i) is the interesting one. Take $y_1 + t_1 a$ and $y_2 + t_2 a$
in $S_a(A)$. We need their distance $\le \diam A$. By the symmetrization
construction, there exist $s_1, s_2$ with $y_1 + s_1 a \in A$ and
$y_2 + s_2 a \in A$, and you can choose them so that
$|s_1 - s_2| \ge |t_1 - t_2|$ (because the original slices are at least
as long as the symmetrized intervals, so you can pick points that are at
least as far apart in the $a$-direction). Then:
\[
  |{(y_1 + t_1 a) - (y_2 + t_2 a)}|^2
  = |y_1 - y_2|^2 + (t_1 - t_2)^2
  \le |y_1 - y_2|^2 + (s_1 - s_2)^2
  \le (\diam A)^2.
\]
Centering the intervals can only bring points in the
$a$-direction closer together. \qedhere
\end{proof}


\subsection{The isodiametric inequality}\label{subsec:isodiametric}

\begin{theorem}[Isodiametric inequality]\label{thm:isodiametric}
For any Borel set $A \subseteq \R^n$,
\[
  \mathcal{L}^n(A) \le \alpha(n)
    \left(\frac{\diam\, A}{2}\right)^{\!n}.
\]
Among all sets of a given diameter, the ball has the largest volume.
\end{theorem}

\begin{proof}
We can assume $\diam A < \infty$ (otherwise the right side is $+\infty$).

\textbf{Step 1: Successive symmetrization.} Let $e_1, \ldots, e_n$ be the
standard basis. Set $A_0 = A$ and $A_k = S_{e_k}(A_{k-1})$ for
$k = 1, \ldots, n$. By Lemma~\ref{lem:steiner-props}:
\[
  \mathcal{L}^n(A_n) = \mathcal{L}^n(A), \qquad
  \diam A_n \le \diam A.
\]
And $A_n$ is symmetric about the origin --- symmetric with respect to each
coordinate hyperplane.

\textbf{Step 2.} Since $A_n$ is origin-symmetric, for every $x \in A_n$
we have $-x \in A_n$. So
\[
  2|x| = |x - (-x)| \le \diam A_n \le \diam A.
\]
This gives $A_n \subseteq B(0, \diam A/2)$.

\textbf{Step 3.}
\[
  \mathcal{L}^n(A) = \mathcal{L}^n(A_n)
  \le \mathcal{L}^n\!\bigl(B(0, \diam A/2)\bigr)
  = \alpha(n)\!\left(\frac{\diam A}{2}\right)^{\!n}. \qedhere
\]
\end{proof}

Symmetrize along each axis. Volume stays the same, diameter can only shrink, and the result is origin-symmetric, so it fits inside a ball of radius $\diam/2$. Done. Three steps, each one line of real content. I love this proof.

\begin{remark}\label{rem:isodiametric-extends}
The theorem extends to Lebesgue measurable sets and even arbitrary sets
(via outer regularity: every set sits inside a Borel set of the same
outer measure).
\end{remark}


\subsection{$\mathcal{H}^n = \mathcal{L}^n$}\label{subsec:Hn-Ln}

This is what all the machinery was building toward.

\begin{theorem}\label{thm:Hn-equals-Ln}
On $\R^n$, $\mathcal{H}^n = \mathcal{L}^n$.
\end{theorem}

\begin{proof}
Two inequalities.

\emph{$\mathcal{H}^n \ge \mathcal{L}^n$:} Let $\{C_j\}$ be any covering
of $A$ with $\diam C_j \le \delta$. By the isodiametric inequality,
\[
  \mathcal{L}^n(A) \le \sum_j \mathcal{L}^n(C_j)
  \le \sum_j \alpha(n)\!\left(\frac{\diam C_j}{2}\right)^{\!n}.
\]
Take the infimum over covers:
$\mathcal{L}^n(A) \le \mathcal{H}^n_\delta(A)$ for all $\delta$, so
$\mathcal{L}^n(A) \le \mathcal{H}^n(A)$.

\emph{$\mathcal{H}^n \le \mathcal{L}^n$:} Cover $A$ by cubes of side
$\delta/\sqrt{n}$ (so diameter $\le \delta$). The sum
$\sum_j \alpha(n)(\diam C_j/2)^n$ is at most $\mathcal{L}^n(A) +
\varepsilon$ if you choose the cover carefully using the definition of
Lebesgue outer measure. Send $\varepsilon \to 0$:
$\mathcal{H}^n_\delta(A) \le \mathcal{L}^n(A)$, so
$\mathcal{H}^n(A) \le \mathcal{L}^n(A)$.

The isodiametric inequality does the heavy lifting in the first direction.
The second direction is checking that the normalization constant works
out. \qedhere
\end{proof}

Brian and I spent a whole afternoon on why $\mathcal{H}^n = \mathcal{L}^n$
is not obvious. You might think: of course covering by balls and covering by
arbitrary sets give the same answer. But they do not, at least not until you prove a
geometric inequality bounding the volume of arbitrary sets by their diameter. That is the isodiametric inequality. Without it you are stuck. We sat in the courtyard outside Kaprielian Hall until it got dark, and I still wasn't sure I believed the second direction until I wrote it up later that night.


\subsection{Densities}\label{subsec:densities}

\begin{definition}[Upper and lower $s$-dimensional densities]%
\label{def:densities}
Let $\mu = \mathcal{H}^s \mres A$. For $x \in \R^n$, define
\begin{align*}
  \Theta^{*s}(\mu, x)
    &= \limsup_{r \to 0^+} \frac{\mu(B(x,r))}{\alpha(s)\, r^s}, \\
  \Theta_*^{\,s}(\mu, x)
    &= \liminf_{r \to 0^+} \frac{\mu(B(x,r))}{\alpha(s)\, r^s}.
\end{align*}
When the limit exists,
$\Theta^s(\mu, x) = \Theta^{*s}(\mu, x) = \Theta_*^{\,s}(\mu, x)$ is
the \emph{$s$-dimensional density} of $\mu$ at $x$.
\end{definition}

Density measures how uniformly $s$-dimensional a set looks at small
scales near $x$. If $A$ is a smooth $k$-manifold, then
$\Theta^k(\mathcal{H}^k \mres A, x) = 1$ at every $x \in A$. Locally
the set looks like a flat $k$-plane, filling the ball at the same rate as
$\R^k$.

\begin{theorem}[Density bounds]\label{thm:density-bounds}
Let $A \subseteq \R^n$ with $\mathcal{H}^s(A) < \infty$. Then:
\begin{enumerate}[nosep]
  \item $\Theta^{*s}(\mathcal{H}^s \mres A,\, x) = 0$ for
        $\mathcal{H}^s$-a.e.\ $x \notin A$.
  \item $2^{-s} \le \Theta^{*s}(\mathcal{H}^s \mres A,\, x) \le 1$ for
        $\mathcal{H}^s$-a.e.\ $x \in A$.
\end{enumerate}
\end{theorem}

The upper bound $\Theta^{*s} \le 1$ is the easier direction. A covering argument using the definition of $\mathcal{H}^s$.
The lower bound $2^{-s} \le \Theta^{*s}$ is harder and relies on a
Vitali-type covering lemma for Hausdorff measure.

\begin{remark}\label{rem:density-integer}
For integer $s = n$ everything is nicer. The $n$-dimensional density of
$\mathcal{L}^n$-measurable sets equals $1$ a.e.\ on the set and $0$
a.e.\ off it. That is the Lebesgue density theorem. For non-integer $s$
the density may not exist, and the gap between $2^{-s}$ and $1$ is
genuine. Sets with non-integer Hausdorff dimension do not have clean local
structure.

I thought this gap was an artifact of the proof until I saw how it shows up in
rectifiability in Chapter~5. It is real.
\end{remark}

TODO: understand the proof of the lower bound $2^{-s}$ better. Evans uses
a covering argument (pp.\ 72--73) that I need to go through line by line.


\subsection{Connection to functions}\label{subsec:functions}

The last bit connects Hausdorff measure back to the ``fine properties''
theme of the book. Main tool: Fubini applied to subgraphs.

\begin{definition}[Subgraph]\label{def:subgraph}
For $f \colon \R^n \to [0, \infty]$ measurable, the \emph{subgraph} is
\[
  \operatorname{SG}(f) = \bigl\{\, (x, t) \in \R^n \times \R :
    0 \le t \le f(x) \,\bigr\} \subseteq \R^{n+1}.
\]
\end{definition}

\begin{lemma}\label{lem:subgraph-measurable}
If $f \colon \R^n \to [0,\infty]$ is $\mathcal{L}^n$-measurable, then
$\operatorname{SG}(f)$ is $\mathcal{L}^{n+1}$-measurable and
\[
  \mathcal{L}^{n+1}\!\bigl(\operatorname{SG}(f)\bigr)
  = \int_{\R^n} f \dd\mathcal{L}^n.
\]
\end{lemma}

\begin{proof}[Proof sketch]
By Fubini:
$\mathcal{L}^{n+1}(\operatorname{SG}(f))
= \int_{\R^n} \mathcal{L}^1(\{t : 0 \le t \le f(x)\})
\dd\mathcal{L}^n(x)
= \int_{\R^n} f(x) \dd\mathcal{L}^n(x)$.
The measurability of $\operatorname{SG}(f)$ as a subset of $\R^{n+1}$
follows from approximation by simple functions. \qedhere
\end{proof}

One of those lemmas that feels obvious but needs Fubini to make honest.
The payoff comes later: you can study integrability, level sets, and
regularity of a function by looking at its subgraph as a subset of
$\R^{n+1}$ and applying Hausdorff measure tools. The Cavalieri principle
($\int f = \int_0^\infty \mathcal{L}^n(\{f > t\})\dd t$) is a close
relative.

\bigskip

That is Chapter~2. Hausdorff measure, Hausdorff dimension, the
isodiametric inequality via Steiner symmetrization, the proof that
$\mathcal{H}^n = \mathcal{L}^n$, densities, and the subgraph
connection. The isodiametric proof is the gem. Everything else is careful bookkeeping with covers and limits.

TODO: read Evans Section~2.4 more carefully for measurability details.
Revisit the Vitali covering argument in the density lower bound.
